\documentclass[11pt,twocolumn,oneside,openany,headings=optiontotoc,11pt,numbers=noenddot]{article}

\usepackage[a4paper]{geometry}
\usepackage[utf8]{inputenc}
\usepackage[T1]{fontenc}
\usepackage{lmodern}
\usepackage[ngerman]{babel}
\usepackage{ngerman}

\usepackage[onehalfspacing]{setspace}

\usepackage{fancyhdr}
\usepackage{fancybox}

\usepackage{rotating}
\usepackage{varwidth}

%Struktogramme
\usepackage[german,curves]{struktex}

\usepackage{pdflscape}
\usepackage{changepage}
\usepackage{graphicx}
\usepackage[bottom]{footmisc}
\usepackage{transparent}
\usepackage{graphbox}
\graphicspath{
	{Pics/PDFs/}
	{Pics/JPGs/}
	{Pics/PNGs/}
}
\usepackage{caption}
\usepackage{wrapfig}
\usepackage{marginnote}
\usepackage{tabularx}
\usepackage{dashrule}
\usepackage{soulutf8}
\usepackage{hhline}
%arydshln suppresses vertical lines in table
%\usepackage{arydshln}
\usepackage{multirow}
\usepackage{enumerate}
\usepackage[hidelinks]{hyperref}
\usepackage{listings}

\usepackage[table]{xcolor}
\usepackage{array}
\usepackage{enumitem,amssymb,amsmath}
\usepackage{interval}
\usepackage{cancel}
\usepackage{stmaryrd}
\usepackage{wasysym}
\usepackage{polynom}
\usepackage{diagbox}
\usepackage{dashrule}
\usepackage{framed}
\usepackage{mdframed}
\usepackage{karnaugh-map}
\usepackage{pdfpages}

\usepackage{blindtext}

\usepackage{eso-pic}

\usepackage{amssymb}
\usepackage{eurosym}

\usepackage[pages=some]{background}
\pagestyle{headings}
\renewcommand{\headrulewidth}{0.2pt}
\renewcommand{\footrulewidth}{0.2pt}
\newcommand*{\underdownarrow}[2]{\ensuremath{\underset{\overset{\Big\downarrow}{#2}}{#1}}}
\setlength{\fboxsep}{5pt}
\newcommand{\explainBelow}[3]{\underbrace{#1}_{\parbox{\widthof{#3}}{\footnotesize\raggedright #2}}}
\newcommand{\explainAbove}[3]{\overbrace{#1}^{\parbox{\widthof{#3}}{\footnotesize\raggedright #2}}}
\newcommand\footnoteref[1]{\protected@xdef\@thefnmark{\ref{#1}}\@footnotemark}


% Codestyle defined
\definecolor{codegreen}{rgb}{0,0.6,0}
\definecolor{codegray}{rgb}{0.5,0.5,0.5}
\definecolor{codepurple}{rgb}{0.58,0,0.82}
\definecolor{backcolour}{rgb}{0.95,0.95,0.92}
\definecolor{deepgreen}{rgb}{0,0.5,0}
\definecolor{darkblue}{rgb}{0,0,0.65}
\definecolor{mauve}{rgb}{0.40, 0.19,0.28}
\colorlet{exceptioncolour}{yellow!50!red}
\colorlet{commandcolour}{blue!60!black}
\colorlet{numpycolour}{blue!60!green}
\colorlet{specmethodcolour}{violet}

%Neue Spaltendefinition
\newcolumntype{L}[1]{>{\raggedright\let\newline\\\arraybackslash\hspace{0pt}}m{#1}}
\newcolumntype{M}{>{\centering\arraybackslash}X}
\newcommand{\cmnt}[1]{\ignorespaces}
%Textausrichtung ändern
\newcommand\tabrotate[1]{\rotatebox{90}{\raggedright#1\hspace{\tabcolsep}}}

%Intervall-Konfig
\intervalconfig {
	soft open fences
}

%Bash
\lstdefinestyle{BashInputStyle}{
	language=bash,
	basicstyle=\small\sffamily,
	backgroundcolor=\color{backcolour},
	columns=fullflexible,
	backgroundcolor=\color{backcolour},
	breaklines=true,
}
%Java
\lstdefinestyle{JavaInputStyle}{
	language=Java,
	backgroundcolor=\color{backcolour},
	aboveskip=1mm,
	belowskip=1mm,
	showstringspaces=false,
	columns=flexible,
	basicstyle={\footnotesize\ttfamily},
	numberstyle={\tiny},
	numbers=none,
	keywordstyle=\color{purple},,
	commentstyle=\color{deepgreen},
	stringstyle=\color{blue},
	emph={out},
	emphstyle=\color{darkblue},
	emph={[2]rand},
	emphstyle=[2]\color{specmethodcolour},
	breaklines=true,
	breakatwhitespace=true,
	tabsize=2,
}
%Python
\lstdefinestyle{PythonInputStyle}{
	language=Python,
	alsoletter={1234567890},
	aboveskip=1ex,
	basicstyle=\footnotesize,
	breaklines=true,
	breakatwhitespace= true,
	backgroundcolor=\color{backcolour},
	commentstyle=\color{red},
	otherkeywords={\ , \}, \{, \&,\|},
	emph={and,break,class,continue,def,yield,del,elif,else,%
		except,exec,finally,for,from,global,if,import,in,%
		lambda,not,or,pass,print,raise,return,try,while,assert},
	emphstyle=\color{exceptioncolour},
	emph={[2]True,False,None,min},
	emphstyle=[2]\color{specmethodcolour},
	emph={[3]object,type,isinstance,copy,deepcopy,zip,enumerate,reversed,list,len,dict,tuple,xrange,append,execfile,real,imag,reduce,str,repr},
	emphstyle=[3]\color{commandcolour},
	emph={[4]ode, fsolve, sqrt, exp, sin, cos, arccos, pi,  array, norm, solve, dot, arange, , isscalar, max, sum, flatten, shape, reshape, find, any, all, abs, plot, linspace, legend, quad, polyval,polyfit, hstack, concatenate,vstack,column_stack,empty,zeros,ones,rand,vander,grid,pcolor,eig,eigs,eigvals,svd,qr,tan,det,logspace,roll,mean,cumsum,cumprod,diff,vectorize,lstsq,cla,eye,xlabel,ylabel,squeeze},
	emphstyle=[4]\color{numpycolour},
	emph={[5]__init__,__add__,__mul__,__div__,__sub__,__call__,__getitem__,__setitem__,__eq__,__ne__,__nonzero__,__rmul__,__radd__,__repr__,__str__,__get__,__truediv__,__pow__,__name__,__future__,__all__},
	emphstyle=[5]\color{specmethodcolour},
	emph={[6]assert,range,yield},
	emphstyle=[6]\color{specmethodcolour}\bfseries,
	emph={[7]Exception,NameError,IndexError,SyntaxError,TypeError,ValueError,OverflowError,ZeroDivisionError,KeyboardInterrupt},
	emphstyle=[7]\color{specmethodcolour}\bfseries,
	emph={[8]taster,send,sendMail,capture,check,noMsg,go,move,switch,humTem,ventilate,buzz},
	emphstyle=[8]\color{blue},
	keywordstyle=\color{blue}\bfseries,
	rulecolor=\color{black!40},
	showstringspaces=false,
	stringstyle=\color{deepgreen}
}

\lstset{literate=%
	{Ö}{{\"O}}1
	{Ä}{{\"A}}1
	{Ü}{{\"U}}1
	{ß}{{\ss}}1
	{ü}{{\"u}}1
	{ä}{{\"a}}1
	{ö}{{\"o}}1
}

% Neue Klassenarbeits-Umgebung
\newenvironment{worksheet}[3]
% Begin-Bereich
{
	\newpage
	\sffamily
	\setcounter{page}{1}
	\ClearShipoutPicture
	\AddToShipoutPicture{
		\put(55,761){{
				\mbox{\parbox{385\unitlength}{\tiny \color{codegray}BBS I Mainz, #1 \newline #2
						\newline #3
					}
				}
			}
		}
		\put(455,761){{
				\mbox{\hspace{0.3cm}\includegraphics[width=0.2\textwidth]{../../logo.pdf}}
			}
		}
	}
}
% End-Bereich
{
	\clearpage
	\ClearShipoutPicture
}

\setlength{\columnsep}{3em}
\setlength{\columnseprule}{0.5pt}

\geometry{left=2.50cm,right=2.50cm,top=3.00cm,bottom=1.00cm,includeheadfoot}
\pagenumbering{gobble}
\pagestyle{empty}

\begin{document}
	\begin{worksheet}{BGY 16}{Klassenstufe 13 - Mathematik}{Lernabschnitt 1: Ableitungsregeln}
		
		\section{Ableitungsregeln}
		Im Verlauf unseres kurzen mathematische Vergnügens werden wir des öfteren mit dem Begriff \textbf{ableiten} oder auch \textbf{Ableitung} konfrontiert.\\
		Nachfolgend werden die verschiedenen Ableitungsregeln kurz erläutert und an einem Beispiel demonstriert.\\
		\par\noindent
		Für die nachfolgenden Regeln betrachten wir die Funktionen \(g\) und \(h\). Diese sind auf einem Intervall \(I\) definiert und an der Stelle \(x\in{}I\) differenzierbar.\\
		Konstanten bezeichnen wir im nachfolgenden mit \(c\in{}\mathbb{R}\).
		\begin{framed}
			\subsection{Potenzregelregel}
			\colorbox{green!10}{\(f(x) = c\cdot{}x^n\)} mit \(c\in{}\mathbb{R}, n\in{}\mathbb{Z}\) ist in \(x\) differenzierbar, dann gilt: \[f'(x) = c\cdot{}n\cdot{}x^{n-1}\]
			\par\noindent
			\underline{Sonderfall}: Steht unser \(x\) im Nenner eines Bruchs, so bietet sich folgendes Vorgehen an:\\
			\(f(x) = \frac{1}{x^3} \xrightarrow{umformen} f(x) = x^{-3}\)\\
			\(\xrightarrow{ableiten} f'(x) = -3\cdot{}x^{\overbrace{-4}^{-3-1}}\)\\
			\par\noindent
			\(\xrightarrow{umformen} f'(x) = \frac{-3}{x^4}\)
			\par\bigskip\noindent
			\rule{0.99\textwidth}{0.4pt}
			\underline{Beispiel:}\[f(x) = x^4\]\[f'(x) = 4x^3\]
		\end{framed}
		\begin{framed}
			\subsection{Summenregel}
			\noindent
			\colorbox{green!10}{\(f(x) = g(x) + h(x)\)} ist in \(x\) differenzierbar, dann gilt: \[f'(x) = g'(x) + h'(x)\]
			\par\noindent
			Als kleine Merkregel führen wir an:\\
			\indent{}Summen werden gliedweise abgeleitet! (Jeder Summand für sich!)
			\par\bigskip\noindent
			\rule{0.99\textwidth}{0.4pt}
			\underline{Beispiel:} \[f(x) = 3x^2 + 4x\]\[f'(x) = 6x + 4\]
		\end{framed}
		\begin{framed}
			\subsection{Faktorregel}
			\colorbox{green!10}{\(f(x) = c\cdot{}g(x)\)} mit \(c\in\mathbb{R}\) ist in \(x\) differenzierbar, dann gilt: \[f'(x) = c\cdot{}g'(x)\]
			\par\noindent
			Unsere Merkregel für die Faktorregel:\\
			Ein konstanter Faktor bleibt beim Ableiten erhalten!
			\par\bigskip\noindent
			\rule{0.99\textwidth}{0.4pt}
			\underline{Beispiel:} \[f(x) = 4x^3\]\[f'(x) = 4\cdot{}3\cdot{}x^2\]
		\end{framed}
		\begin{framed}
			\subsection{Kettenregel}
			Betrachten wir \colorbox{green!10}{\(f(x) = g(x)\circ{}h(x) = g(h(x))\)} und sowohl \(g(x)\) als auch \(h(x)\) sind differenzierbar in \(x\), so bezeichnet man \(f(x)\) auch als Verkettung der beiden Funktionen \(g\) und \(h\). Dann gilt die folgende Ableitungsregel: \[f'(x) = g'(h(x))\cdot{}h'(x)\]
			\par\noindent
			Im Allgemeinen können wir diese Regel auch formulieren als \textit{äußere Ableitung} mal \textit{innere Ableitung}.\\
			\par\bigskip\noindent
			\rule{0.99\textwidth}{0.4pt}
			\underline{Beispiel:} \[f(x) = \underbrace{(\overbrace{x^2-3x}^{h(x)})^4}_{g(x)}\]\[f'(x) = \underbrace{4(x^2-3x)^3}_{g'(h(x))}\cdot{}\underbrace{(2x-3)}_{h'(x)}\]
		\end{framed}
		\begin{framed}
			\subsection{Produktregel}
			\colorbox{green!10}{\(f(x) = g(x)\cdot{}h(x)\)} ist in \(x\) differenzierbar, dann gilt: \[f'(x) = g'(x)\cdot{}h(x) + g(x)\cdot{}h'(x)\]
			\par\noindent
			Prinzipiell können wir uns für diese Ableitung auch schnell merken: (uv)' = u'v + uv'
			\par\bigskip\noindent
			\rule{0.99\textwidth}{0.4pt}
			\underline{Beispiel:}\[f(x) = \underbrace{x}_{g(x)}\cdot{}\underbrace{(x^3+x)}_{h(x)}\]\[f'(x) = \underbrace{1}_{g'(x)}\cdot\underbrace{(x^3+x)}_{h(x)} + \underbrace{x}_{g(x)}\cdot\underbrace{(3x^2+1)}_{h'(x)}\]
		\end{framed}
		\begin{framed}
			\subsection{Quotientenregel}
			Sind \(g(x)\) und \(h(x)\) differenzierbar und \underline{\(h(x)\neq{}0\)}, so ist auch \colorbox{green!10}{\(f(x) = \frac{g(x)}{h(x)}\)} differenzierbar. Die entsprechende Ableitungsregel lautet wie folgt: \[f'(x) = \frac{g'(x)\cdot{}h(x) - g(x)\cdot{}h'(x)}{h^2(x)}\]
			\par\noindent
			Prinzipiell können wir uns für diese Ableitung auch schnell merken: (uv)' = u'v + uv'
			\par\bigskip\noindent
			\rule{0.99\textwidth}{0.4pt}
			\underline{Beispiel:}\[f(x) = \frac{\overbrace{(x^3+2x)}^{g(x)}}{\underbrace{x^2}_{h(x)}}\]\[f'(x) = \frac{\overbrace{(3x^2+2)}^{g'(x)}\cdot\overbrace{x^2}^{h(x)} - \overbrace{(x^3+2x)}^{g(x)}\cdot\overbrace{2x}^{h'(x)}}{\underbrace{(x^2)^2}_{h(x)}}\]
		\end{framed}
	\end{worksheet}
\end{document}